\section{Research program and falsifiable next steps}
\label{sec:research}

The framework's contribution at this stage is a precise separation of
semantics, compilation contracts, and evidence types, with elementary
preservation results. It does not yet establish a new separation theorem or
a complete basis of all nonlocal games. The following questions would turn
the architecture into a substantive research program.

\paragraph{A useful basis of constructors.}
Express benchmark families using parity, consistency, sampling, local wiring,
hidden/public selection, and parallel composition. Measure description size,
lowered size, compilation cost, and analysis cost separately. Investigate
which additional constructors are needed for linear-system, graph-coloring,
and synchronous games. Universality through arbitrary lookup tables is easy;
a basis that preserves exploitable structure is the meaningful target.

\paragraph{Certificates that compose.}
Start with component value bounds and SOS relations. Use the deficit rule of
Section~\ref{sec:composition} to track robustness under hidden mixtures, then
identify sufficient compatibility conditions for the component extraction
maps. A theorem should quantify the cost in score deficit, local dimension,
and target distance. Merely adding two self-testing scores is not a proof
that their certified subsystems coexist as a tensor product.

\paragraph{Resource-aware optimization.}
Construct pairs of arenas with identical payoff syntax but distinct source,
dimension, or trust assumptions. Ensure that the checker rejects invalid
backend reuse and that every accepted relaxation has the documented bound
direction. The no-input obstruction in Proposition~\ref{prop:noinput} is a
particularly useful test: a network compatibility problem should not be
mistaken for an ordinary linear game-value separation.

\paragraph{Certified numerical discovery.}
Search over structured game parameters, producing a candidate strategy lower
bound and a local upper bound for separation. Then seek a quantum upper bound
and reconstruct an exact dual or SOS certificate. Keep candidate discovery
and proof checking separate. Optimizing gaps over unconstrained score tables
is ill-posed because scaling increases the gap; normalize sampling/payoff and
control alphabet size, coefficient complexity, and physical noise assumptions.

\paragraph{A staged benchmark suite.}
CHSH tests value conventions, XOR specialization, an explicit witness, and
an exact SOS identity. GHZ parity tests expose multipartite semantics.
Magic-square consistency tests exercise nonbinary answer alphabets and
perfect quantum strategies. Degenerate zero-weight questions, biased
sampling, negative general payoffs, and no-input games test boundary cases.
Subsequent stages should add a dimension witness, a nonlinear network
compatibility problem, a semiquantum game, and a genuinely sequential
protocol. Each addition should come with an explicit semantic extension and
a demonstrably valid analysis method.

\paragraph{Open design question.}
Can a modest typed causal IR support all of these fragments while keeping
resource-preservation proofs small and backend-specific? This is a research
question about a useful representation, not a conjecture that every game
admits efficient solution. Success can be measured by reuse, verified
lowerings, smaller descriptions, and reproducible certificates even when
many values remain unresolved.
